\documentclass[10pt,a4paper]{amsart}
\usepackage[margin=19mm]{geometry}
\usepackage{amsmath,amssymb,mathtools}
\usepackage[scr=rsfs,cal=euler]{mathalfa}
\usepackage{xfrac}
\usepackage[unicode,hidelinks,bookmarks=false]{hyperref}
\newcommand{\ie}{i.e.\ }
\newcommand{\cf}{cf.\ }
\newcommand{\resp}{respectively\ }
\newcommand{\Subsection}{\subsection}
\providecommand{\nobreakdash}{\nobreak\hbox{-}}
\setlength{\emergencystretch}{2em}
\multlinegap0pt
\usepackage{bm}
\usepackage{bbm}
\usepackage{dsfont}
\usepackage{xspace}
\usepackage{cancel}
\usepackage{mathtools}
\usepackage{amsthm}
\makeatletter
\@ifundefined{thm}{\newtheorem{thm}{Thm}}{}
\@ifundefined{lem}{\newtheorem{lem}[thm]{Lemma}}{}
\@ifundefined{prop}{\newtheorem{prop}[thm]{Proposition}}{}
\makeatother
\let\ol=\overline
\newcommand{\Aut}{{\rm Aut}}
\newcommand{\Id}{{\rm Id}}
\newcommand{\N}{{\mathbb N}}
\newcommand{\R}{{\mathbb R}}
\newcommand{\Sc}{{\mathbb S}}
\newcommand{\Sub}{{\rm Sub}}
\newcommand{\du}{{\rm d\,}}
\newcommand{\lf}{{\mathfrak l}}
\newcommand{\spg}{{\rm Sub}^p_G}
\newcommand{\spp}{{\rm Sub}^{p}}
\title[Characterisation of distal actions of automorphisms on the s]{Characterisation of distal actions of automorphisms on the space of one-parameter subgroups of Lie groups}
\author{Debamita Chatterjee, Riddhi Shah}
\date{Source: arXiv:2406.01237v3 (CC BY 4.0)}
\begin{document}
\maketitle

\noindent\emph{Source and licence.} Characterisation of distal actions of automorphisms on the space of one-parameter subgroups of Lie groups. Version arXiv:2406.01237v3 (CC BY 4.0), distributed under the Creative Commons Attribution 4.0 International licence, as recorded in the arXiv licence metadata for this exact version and shown on the abstract page https://arxiv.org/abs/2406.01237v3. Full rights notice: licenses/arxiv-2406.01237-CC-BY-4.0.txt.\par\smallskip

\noindent\textit{Editorial scope.} Counted unit: the Prop whose source label is \texttt{crux} in the article (source0.tex lines 799--802, source label \texttt{crux}), delivered together with the article's own complete original proof of it (source0.tex lines 804--875, one \texttt{proof} environment of 72 lines). No mathematics is written, repaired or re-derived: statement and proof are the article's own text line for line. Required context retained uncounted: discrete (lines 387--390); gk (lines 613--617). Every definition, display and label of those lines is retained unchanged. Omitted from this package and not redistributed: 1--386, 391--612, 618--798, 803--803, 876--1277. Nothing inside a retained excerpt is omitted or altered: every retained body is delivered exactly as the article's lines are written, so no omission receipt of any kind accompanies this package. Citations: the retained excerpts carry no citation command at all, so no bibliography entry of the article is transcribed and no reference is redistributed. Reference adaptation: none was needed and none was made; every reference inside the delivered unit resolves inside it, so no unresolved reference or citation is produced. Printed slips kept literal: no printed word, symbol, hypothesis or number of the counted statement or of its proof was altered. Layout and class: the article's own class is \texttt{amsart} with packages including graphicx, amssymb, amsmath, setspace, color, hyperref, cleveref, xcolor, xfrac, amsfonts, bm; the wrapper is the lane's pinned \texttt{amsart} editorial wrapper at 10pt on A4, which restates the theorem-like environments the retained text needs and adds the article's own macro definitions that the retained text needs (\texttt{\textbackslash Aut}, \texttt{\textbackslash Id}, \texttt{\textbackslash N}, \texttt{\textbackslash R}, \texttt{\textbackslash Sc}, \texttt{\textbackslash Sub}, \texttt{\textbackslash du}, \texttt{\textbackslash lf}, \texttt{\textbackslash ol}, \texttt{\textbackslash spg}, \texttt{\textbackslash spp}), with extra wrapper packages (bm, bbm, dsfont, xspace, cancel, mathtools). The delivered numbering is the unit's own. Assets: no retained span contains a figure environment, a graphics-inclusion command, a PDF-page inclusion, a table or a TikZ picture; no image file is copied into this package, the package declares no asset dependency and no image file is redistributed with it. Licence: version arXiv:2406.01237v3 of this article is distributed under the Creative Commons Attribution 4.0 International licence, as recorded in the arXiv licence metadata for this exact version and shown on the abstract page https://arxiv.org/abs/2406.01237v3. A full rights notice is delivered at licenses/arxiv-2406.01237-CC-BY-4.0.txt. Characters: every retained excerpt is pure ASCII, so the delivered engine, XeLaTeX with its default Latin Modern text font, reports no missing character.
\begin{lem} \label{discrete}
Let $G$ be a connected Lie group.  Then the trivial subgroup $\{e\}$ is isolated in $\Sub^{co}_G$, and so in $\spg$. Moreover, 
$\Sub^{co}_G$ does not contain any non-trivial discrete subgroup.
\end{lem} 

\begin{prop} \label{gk}
Let $G$ be a connected Lie group and let $T\in\Aut(G)$. Let $K$ be a $T$-invariant compact connected 
central subgroup of $G$ and let $\ol{T}$ denote the automorphism of $G/K$ induced by $T$. 
If $T$ acts distally on $\spg$, then $\ol{T}$ acts distally on $\spp_{G/K}$. 
\end{prop}

\begin{prop} \label {crux}
Let $G$ be a connected Lie group, $T\in\Aut(G)$ and let $L$ be a closed connected nilpotent normal $T$-invariant 
subgroup of $G$. Suppose $T$ acts trivially on both $L$ and $G/L$. If $T$ acts distally on $\spg$, then $T=\Id$. 
\end{prop}

\begin{proof} Since $T$ acts trivially on $G/L$ and $L$, we have that $T$ acts distally on $G$. In fact, all the eigenvalues 
of $\du T$ are equal to 1 and $T$ is unipotent. Suppose $T$ acts distally on $\spg$. Let $U$ be a neighbourhood 
of the identity $e$ such that every element of $U$ is contained in a one-parameter subgroup. Since $G$ is 
connected, $U$ generates $G$ and it is enough to show that $T(x)=x$ for all $x\in U$. 
 
\smallskip
 \noindent{\bf Step 1:} If possible, suppose $T(x)\ne x$ for some $x\in U$. There exists a one-parameter 
 subgroup $\{x_t\}$ in $G$ such that $x_1=x$. Let $H=\ol{\{x_t\}L}$. Then $H$ is a closed connected solvable $T$-invariant 
 subgroup of $G$ and $H/L$ is abelian. Let $\pi:G\to G/L$ be the natural projection.  
 Then $\pi(H)=\ol{\pi(\{x_t\})}$. If $\pi(\{x_t\})$ is unbounded, it is closed and it is isomorphic to $\R$ and in this case 
 $\{x_t\}$ is also unbounded and closed. Moreover $H=\{x_t\}L$.
 
Now suppose $\pi(\{x_t\})$ is relatively compact. As $H$ is solvable, there exists a maximal compact 
(connected, abelian) subgroup $K$ in $H$ such that $H=KL$. As $T(x)\ne x$ and $T|_L=\Id$, 
we have that $T(h)\ne h$ for some $h\in K$, and hence $T(h)\ne h$ for some finite order element $h$ in $K\cap U$.
There exists a closed (compact) one-parameter subgroup $\{h_t\}$ in $K$ such that $h_1=h$. Replacing 
$\{x_t\}$ by $\{h_t\}$, we have in this case too that $H$ is closed and $H=\{x_t\}L$.

Now we have that $H=\{x_t\}L$ is closed, where either $\pi(\{x_t\})$ (as well as $\{x_t\}$) is isomorphic to $\R$ 
or $\{x_t\}$ is isomorphic to $\Sc^1$ and $T(x_1)\ne x_1$. As $T$ acts trivially on $G/L$, we have that $T(H)=H$ 
and $T(x_t)=x_ty_t$, for some $y_t\in L$, $t\in\R$, where $y_1\ne e$. 

 \smallskip
 \noindent{\bf Step 2:} Let $A:=\{x_t\}$ and let $\{n_k\}\subset\N$ be an unbounded sequence 
such that $T^{n_k}(A)\to B$ for some $B\in\spp_H$. By Lemma \ref{discrete}, $B$ is not discrete as $A$ 
is connected and non-trivial. Let $b\in B$ be such that $b\ne e$. There exist $t_k\in\R$, $k\in\N$, 
such that 
$$
T^{n_k}(x_{t_k})=x_{t_k}y_{t_k}^{n_k}\to b.$$ 
 As $y_t\in L$, $t\in\R$, we get that $\pi(x_{t_k})\to\pi(b)$. If $H/L$ is isomorphic to $\R$, 
then we get that the sequence $\{t_k\}$ is bounded. In the second case where $\{x_t\}$ is compact, we 
can choose $\{t_k\}$ to be bounded. Passing to a subsequence if necessary, we may assume that 
$t_k\to t_0$. Therefore, $x_{t_k}\to x_{t_0}$ and $y_{t_k}\to y_{t_0}$ and 
$b=x_{t_0}y'$, where $y_{t_k}^{n_k}\to y'$ in $L$. 

 \smallskip
 \noindent{\bf Step 3:} We first assume that $L$ is simply connected. We show that $y_{t_0}=e$ in this case.
Let $\lf$ be the Lie algebra of $L$. Since $L$ is simply connected and nilpotent, $\exp: \lf\to L$ is a 
homeomorphism with $\log$ as its inverse. Let $y_{t_k}=\exp Y_k$ for some $Y_k\in\lf$. 
Then $y_{t_k}^{n_k}=\exp(n_kY_k)$ and as $y_{t_k}^{n_k}\to y'$, we have 
that $n_kY_k\to \log y'$ in $\lf$. As $\lf$ is isomorphic to $\R^d$ for some $d\in\N$, it follows that $Y_k\to 0$ in $\lf$, 
and hence $y_{t_k}\to e$, i.e.\ $y_{t_0}=e$. Now
$$
T(b)=T(x_{t_0})y'=x_{t_0}y_{t_0}y'=x_{t_0}y'=b.$$
Thus $T(B)=B$. As $T$ acts distally on $\spg$, we have that $T(A)=A$. Since $A$ is a one-parameter subgroup in $G$ 
and $T$ is unipotent, we get that $T(x_t)=x_t$, $t\in\R$. This leads to a contradiction as $T(x_1)\ne x_1$. 
Therefore, $T(x)=x$ for all $x\in U$, and hence, for all $x\in G$. Thus $T=\Id$ if $L$ is simply connected. 
 
 \smallskip
 \noindent{\bf Step 4:} Suppose that $L$ is not simply connected. Let $M$ be the maximal 
 compact (normal) subgroup of $L$. Then $M$ is central in $G$ and $L/M$ is simply connected. 
 Let $\varrho:G\to G/M$ be the natural projection and let $\ol{T}$ be the automorphism on $G/M$ 
 induced by $T$. Then $\ol{T}$ acts trivially on both $\varrho(L)$ and $\varrho(G)/\varrho(L)$. 
 By Proposition \ref{gk}, $\ol{T}$ acts distally on $\spp_{G/M}$. From Step 3, we get that $\ol{T}$ acts trivially on 
 $G/M$, i.e.\ $T(g)\in gM$, $g\in G$. In particular, for $\{x_t\}$ as above, $T(x_t)=x_ty_t$, where $y_t\in M$, 
 for all $t\in\R$, and $y_1\ne e$. Since $M$ is central in $G$, $y_tx_s=x_sy_t$, $t, s\in \R$, and $\{y_t\}$ is a 
 one-parameter subgroup of the compact central subgroup $M$. Note that $\{y_t\}$ is closed if $\{x_t\}$ is compact. 

As $\spg$ is compact, there exists an unbounded sequence $\{n_k\}$ in $\N$ such that $T^{n_k}(\{x_t\})\to C$ 
for some $C\in\spg$. We show that $C=\{x_t\}\times \ol{\{y_t\}}$. As $\{x_t\}\subset \{x_t\}\times \ol{\{y_t\}}$, 
and the latter is $T$-invariant, we get that $C\subset \{x_t\}\times \ol{\{y_t\}}$. For a fixed $t\in\R$, 
$$
T^{n_k}(x_{t/n_k})=x_{t/n_k}y_{t/n_k}^{n_k}=x_{t/n_k}y_t\to y_t$$ 
as $n_k\to\infty$. Therefore, $y_t\in C$ for each $t$, and $\ol{\{y_t\}}\subset C$ as $C$ is closed. Also, 
$T^{n_k}(x_t)=x_ty_t^{n_k}\to x_ty$ along a subsequence of $\{n_k\}$ for some $y\in\ol{\{y_t\}}\subset C$.
Hence, $x_t\in C$, $t\in\R$. Thus $\{x_t\}\times \ol{\{y_t\}}\subset C$. 
Now $T(C)=C$, and we get that $T^{n_k}(T(\{x_t\}))\to C$. As $T$ acts distally on $\spg$, we have that 
$T(\{x_t\})=\{x_t\}$, and since $T$ is unipotent, $T(x_t)=x_t$ for each $t$. 
But $T(x_1)=x_1y_1\ne x_1$.  
This leads to a contradiction. Therefore, $T(x)=x$ for each $x\in U$, and hence, for each $x\in G$. 
Thus $T=\Id$.
 \end{proof}

\end{document}
